% GENERATED FILE. Edit canonical lessons, then run npm run generate:books.
% canonical-source-hash: fnv1a64:4f67c864c81b2a46
% canonical-lessons: UR-C388-signal, UR-C388-jankshan, UR-C388-phatak, UR-C388-barth, UR-C388-takna

\chapter{Signals and Crossings}
\label{ch:ur-signals-and-crossings}
\hlchaptermodality{388}

\begin{chapteropening}
\textbf{By the end of this chapter:} \emph{I can say what the signal shows and why the crossing is shut.}

Complete the last lesson of chapter 388: I can say what the signal shows and why the crossing is shut.
\end{chapteropening}

\section[\texorpdfstring{\ur{سگنل} (signal)}{سگنل (signal)}]{\ur{سگنل} (signal) — a signal (for trains or traffic)}

\label{lesson:UR-C388-signal}



\begin{quote}
Before the new one: say the Urdu for to haul, then the Urdu for to seat (someone).
\end{quote}

\subsection*{You'll want to know: \ur{سگنل}}
\textbf{\ur{سگنل}} — \emph{signal} — \textquotedblleft{}a signal (for trains or traffic)\textquotedblright{}.

\begin{quote}
\ur{سگنل} \ur{لال} \ur{تھا،} \ur{اس} \ur{لیے} \ur{گاڑی} \ur{رک} \ur{گئی۔}
\end{quote}

\begin{quote}
\emph{signal lāl thā, is liye gāṛī ruk gaī.}
\end{quote}

\begin{quote}
The signal was red, so the train stopped.
\end{quote}

Masculine.

\begin{cousinweb}
\textbf{\ur{سگنل}} (\emph{signal}) is borrowed from English \emph{signal}.
\end{cousinweb}

\begin{grammarlens}[title={What you've built}]
One more word for getting around.
\end{grammarlens}

\subsection*{Your turn}
\begin{itemize}
\raggedright
  \item \emph{Say it:} \emph{signal}
  \item \emph{Say it:} \emph{signal}, once more
  \item \emph{Say it:} \emph{signal lāl thā, is liye gāṛī ruk gaī.} — The signal was red, so the train stopped.
  \item \emph{Recall:} say the Urdu for to haul, then the Urdu for to seat (someone), then say \emph{signal} again
\end{itemize}

\begin{tcolorbox}[breakable,colback=teal!4,colframe=teal!35!black,title={Before you move on}]
What does \ur{سگنل} mean? (\textquotedblleft{}A signal (for trains or traffic)\textquotedblright{}.) Say it once more.
\end{tcolorbox}
\section[\texorpdfstring{\ur{جنکشن} (jankshan)}{جنکشن (jankshan)}]{\ur{جنکشن} (jankshan) — a railway junction}

\label{lesson:UR-C388-jankshan}



\begin{quote}
Before the new one: say the Urdu for to seat (someone), then the Urdu for a signal (for trains or traffic).
\end{quote}

\subsection*{You'll want to know: \ur{جنکشن}}
\textbf{\ur{جنکشن}} — \emph{jankshan} — \textquotedblleft{}a railway junction\textquotedblright{}.

\begin{quote}
\ur{یہ} \ur{اسٹیشن} \ur{ایک} \ur{بڑا} \ur{جنکشن} \ur{ہے۔}
\end{quote}

\begin{quote}
\emph{yih isṭeshan ek baṛā jankshan hai.}
\end{quote}

\begin{quote}
This station is a big junction.
\end{quote}

Masculine.

\begin{cousinweb}
\textbf{\ur{جنکشن}} (\emph{jankshan}) is borrowed from English \emph{junction}: a station where railway lines meet.
\end{cousinweb}

\begin{grammarlens}[title={What you've built}]
One more word for getting around.
\end{grammarlens}

\subsection*{Your turn}
\begin{itemize}
\raggedright
  \item \emph{Say it:} \emph{jankshan}
  \item \emph{Say it:} \emph{jankshan}, once more
  \item \emph{Say it:} \emph{yih isṭeshan ek baṛā jankshan hai.} — This station is a big junction.
  \item \emph{Recall:} say the Urdu for to seat (someone), then the Urdu for a signal (for trains or traffic), then say \emph{jankshan} again
\end{itemize}

\begin{tcolorbox}[breakable,colback=teal!4,colframe=teal!35!black,title={Before you move on}]
What does \ur{جنکشن} mean? (\textquotedblleft{}A railway junction\textquotedblright{}.) Say it once more.
\end{tcolorbox}
\section[\texorpdfstring{\ur{پھاٹک} (phāṭak)}{پھاٹک (phāṭak)}]{\ur{پھاٹک} (phāṭak) — a gate; a level crossing}

\label{lesson:UR-C388-phatak}



\begin{quote}
Before the new one: say the Urdu for a signal (for trains or traffic), then the Urdu for a railway junction.
\end{quote}

\subsection*{You'll want to know: \ur{پھاٹک}}
\textbf{\ur{پھاٹک}} — \emph{phāṭak} — \textquotedblleft{}a gate; a level crossing\textquotedblright{}.

\begin{quote}
\ur{گاڑی} \ur{آ} \ur{رہی} \ur{تھی،} \ur{اس} \ur{لیے} \ur{پھاٹک} \ur{بند} \ur{تھا۔}
\end{quote}

\begin{quote}
\emph{gāṛī ā rahī thī, is liye phāṭak band thā.}
\end{quote}

\begin{quote}
A train was coming, so the level crossing was shut.
\end{quote}

Masculine.

\begin{cousinweb}
\textbf{\ur{پھاٹک}} (\emph{phāṭak}) is inherited from \emph{phāṭṭakka}, \textquotedblleft{}a gate\textquotedblright{}, a form CDIAL reconstructs (CDIAL 9069).
\end{cousinweb}

\begin{grammarlens}[title={What you've built}]
One more word for getting around.
\end{grammarlens}

\subsection*{Your turn}
\begin{itemize}
\raggedright
  \item \emph{Say it:} \emph{phāṭak}
  \item \emph{Say it:} \emph{phāṭak}, once more
  \item \emph{Say it:} \emph{gāṛī ā rahī thī, is liye phāṭak band thā.} — A train was coming, so the level crossing was shut.
  \item \emph{Recall:} say the Urdu for a signal (for trains or traffic), then the Urdu for a railway junction, then say \emph{phāṭak} again
\end{itemize}

\begin{tcolorbox}[breakable,colback=teal!4,colframe=teal!35!black,title={Before you move on}]
What does \ur{پھاٹک} mean? (\textquotedblleft{}A gate; a level crossing\textquotedblright{}.) Say it once more.
\end{tcolorbox}
\section[\texorpdfstring{\ur{برتھ} (barth)}{برتھ (barth)}]{\ur{برتھ} (barth) — a berth (on a train)}

\label{lesson:UR-C388-barth}



\begin{quote}
Before the new one: say the Urdu for a railway junction, then the Urdu for a gate.
\end{quote}

\subsection*{You'll want to know: \ur{برتھ}}
\textbf{\ur{برتھ}} — \emph{barth} — \textquotedblleft{}a berth (on a train)\textquotedblright{}.

\begin{quote}
\ur{میری} \ur{برتھ} \ur{اوپر} \ur{ہے۔}
\end{quote}

\begin{quote}
\emph{merī barth ūpar hai.}
\end{quote}

\begin{quote}
My berth is the top one.
\end{quote}

Feminine.

\begin{cousinweb}
\textbf{\ur{برتھ}} (\emph{barth}) is borrowed from English \emph{berth}: a bed on a night train.
\end{cousinweb}

\begin{grammarlens}[title={What you've built}]
One more word for getting around.
\end{grammarlens}

\subsection*{Your turn}
\begin{itemize}
\raggedright
  \item \emph{Say it:} \emph{barth}
  \item \emph{Say it:} \emph{barth}, once more
  \item \emph{Say it:} \emph{merī barth ūpar hai.} — My berth is the top one.
  \item \emph{Recall:} say the Urdu for a railway junction, then the Urdu for a gate, then say \emph{barth} again
\end{itemize}

\begin{tcolorbox}[breakable,colback=teal!4,colframe=teal!35!black,title={Before you move on}]
What does \ur{برتھ} mean? (\textquotedblleft{}A berth (on a train)\textquotedblright{}.) Say it once more.
\end{tcolorbox}
\section[\texorpdfstring{\ur{تاکنا} (tāknā)}{تاکنا (tāknā)}]{\ur{تاکنا} (tāknā) — to gaze (at); to look out (at)}

\label{lesson:UR-C388-takna}



\begin{quote}
Before the new one: say the Urdu for a gate, then the Urdu for a berth (on a train).
\end{quote}

\subsection*{You'll want to know: \ur{تاکنا}}
\textbf{\ur{تاکنا}} — \emph{tāknā} — \textquotedblleft{}to gaze (at); to look out (at)\textquotedblright{}.

\begin{quote}
\ur{سارا} \ur{کھڑکی} \ur{سے} \ur{کھیتوں} \ur{کو} \ur{تاکتی} \ur{ہے۔}
\end{quote}

\begin{quote}
\emph{sārā khiṛkī se khetoṅ ko tāktī hai.}
\end{quote}

\begin{quote}
Sara gazes out of the window at the fields.
\end{quote}

Past, with \ur{نے} (he or she …): \ur{اس} \ur{نے} \ur{کھڑکی} \ur{سے} \ur{باہر} \ur{تاکا}, \emph{us ne khiṛkī se bāhar tākā}.

\begin{cousinweb}
\textbf{\ur{تاکنا}} (\emph{tāknā}) is inherited: CDIAL lists the Hindi cousin \emph{tāknā}, \textquotedblleft{}to look for\textquotedblright{}, under Sanskrit \emph{tarkayati}, \textquotedblleft{}guesses, thinks\textquotedblright{} (CDIAL 5716); the sense \textquotedblleft{}to gaze at\textquotedblright{} probably grew from looking out for something.
\end{cousinweb}

\begin{grammarlens}[title={What you've built}]
One more word for getting around.
\end{grammarlens}

\subsection*{Your turn}
\begin{itemize}
\raggedright
  \item \emph{Say it:} \emph{tāknā}
  \item \emph{Say it:} \emph{tāknā}, once more
  \item \emph{Say it:} \emph{sārā khiṛkī se khetoṅ ko tāktī hai.} — Sara gazes out of the window at the fields.
  \item \emph{Recall:} say the Urdu for a gate, then the Urdu for a berth (on a train), then say \emph{tāknā} again
\end{itemize}

\begin{tcolorbox}[breakable,colback=teal!4,colframe=teal!35!black,title={Before you move on}]
What does \ur{تاکنا} mean? (\textquotedblleft{}To gaze (at); to look out (at)\textquotedblright{}.) Say it once more.
\end{tcolorbox}
